%=== spec/CA/calculalgeb_12.tex || CA || Algébrique

Lorsqu'un point mobile suit une trajectoire circulaire de rayon $R$, en mètre (m), son accélération centripète $a$ (en $\text{m}/\text{s}^{2}$) s'exprime en fonction de la vitesse $v$ (en $\text{m}/\text{s}$) de la manière suivante : \[
a=\frac{v^{2}}{R}. \]

L'expression permettant, à partir de cette formule, d'exprimer la vitesse $v$ est :

\smallskip

\eamreponsesautom[2]{$v=a R^{2}$}{$v=\sqrt{a R}$}{$v=\sqrt{\frac{a}{R}}$}{$v=\frac{a^{2}}{R}$}